\chapter{14. Planificación académica general}
\label{chap:planificacion-academica}

\section{Criterio de planificación}

La planificación conserva el calendario académico comprendido entre marzo y diciembre de 2026, pero diferencia dos niveles. El primero reúne los hitos formales de la asignatura: propuesta, entregas del 25 \%, 50 \% y 75 \%, exposición y defensa. El segundo organiza la evolución técnica del producto mediante checkpoints, dependencias entre repositorios y evidencia de cierre.

Esta separación evita medir el avance únicamente por cantidad de código. Cada bloque se considera cerrado cuando cuenta con implementación o documentación, criterio de aceptación, evidencia y conexión con el roadmap general.

\section{Planificación por entregables}

% TABLE_PENDING_MIGRATION
% DOCX_TABLE_NUMBER: 14.1
% DOCX_TABLE_TITLE: Roadmap actualizado para la entrega del 50 %.
% TABLE_ROLE: DELIVERABLE_PLAN
% TABLE_ROWS: 18
% TABLE_COLUMNS: 3
% TABLE_MERGED_CELLS: NO
% TABLE_COMPLEXITY: COMPLEX

\section{Gestión de cambios}

Los cambios de alcance deben registrarse con su origen. La incorporación axial se vincula con user research y factibilidad de datos. El visor volumétrico se vincula con la constatación de que los MHA ya contienen múltiples cortes y con la necesidad de aproximar la experiencia a un lector real. La comparación longitudinal surge como evolución de producto, pero permanece fuera del alcance inmediato.

Cada cambio debe actualizar requerimientos, roadmap, arquitectura, contratos, pruebas y documentación. No corresponde modificar solo la interfaz cuando la capacidad depende de nuevos assets o persistencia.

\section{Riesgos de planificación}

Los principales riesgos son: sobrealcance clínico, apertura de funcionalidades antes de cerrar dependencias, falta de evidencia, divergencia entre repositorios, costos cloud no controlados, ausencia de validación profesional sobre la versión integrada e inconsistencias entre documentación y producto.

Las mitigaciones son: checkpoints compartidos, commits coordinados, criterios de aceptación, separación entre implementado y planificado, control presupuestario, pruebas por capas y uso del documento maestro como fuente editorial común.

\section{Conclusión de la planificación}

La planificación no reinicia el proyecto después del 25 \%: toma como base los modelos, la integración multiplanar y la arquitectura ya desarrollados. El camino crítico ya no pasa por generar y persistir las vistas previas de todos los cortes, capacidad hoy resuelta, sino por integrar en el frontend los contratos congelados en P10.9, cerrar el recorrido extremo a extremo con interfaz, realizar la validación profesional y avanzar hacia el despliegue. La región de interés axial automática constituye una línea paralela que no condiciona ninguno de esos pasos.
